
\documentclass{article}
\usepackage{colortbl}
\usepackage{makecell}
\usepackage{multirow}
\usepackage{supertabular}

\begin{document}

\newcounter{utterance}

\centering \large Interaction Transcript for game `wordle', experiment `high\_frequency\_words\_no\_clue\_no\_critic', episode 8 with gpt{-}oss{-}20b{-}t1.0.
\vspace{24pt}

{ \footnotesize  \setcounter{utterance}{1}
\setlength{\tabcolsep}{0pt}
\begin{supertabular}{c@{$\;$}|p{.15\linewidth}@{}p{.15\linewidth}p{.15\linewidth}p{.15\linewidth}p{.15\linewidth}p{.15\linewidth}}
    \# & \multicolumn{2}{c}{Player} && \multicolumn{2}{c}{Game Master} \\
    \hline

    \theutterance \stepcounter{utterance}  
    & & & \multicolumn{4}{p{0.6\linewidth}}{
        \cellcolor[rgb]{0.9,0.9,0.9}{
            \makecell[{{p{\linewidth}}}]{
                \texttt{\tiny{[P1$\langle$GM]}}
                \texttt{You are a language wizard who likes to guess words by using the given rules.} \\
\\ 
\texttt{Welcome to Wordle! You have six attempts to guess the target word, a valid English word of five lowercase letters (a{-}z). Please use the tags "explanation:" and "guess:" to provide a concise explanation for each guess.} \\
\\ 
\texttt{For instance, if your guess is "apple", your response should be} \\
\texttt{explanation: this is a common five{-}letter English word, and I am starting my guess with this word.} \\
\texttt{guess: apple} \\
\\ 
\texttt{After each guess, your answer will be validated, and you will receive feedback indicating which letters are correct (green), which letters are correct but in the wrong position (yellow), and which letters are incorrect (red). This feedback can be useful in determining which letters to include or exclude in your next guess.} \\
\\ 
\texttt{For example, the feedback for "apple" might be:} \\
\texttt{guess\_feedback: a<yellow> p<yellow> p<green> l<yellow> e<red>} \\
\\ 
\texttt{The explanation should contain details about how the guess\_feedback is used to arrive at a new guess.} \\
\\ 
\texttt{Do not generate any other text.} \\
\\ 
\texttt{Let's begin with your first guess.} \\
            }
        }
    }
     \\ \\

    \theutterance \stepcounter{utterance}  
    & \multicolumn{4}{p{0.6\linewidth}}{
        \cellcolor[rgb]{0.561,0.737,0.561}{
            \makecell[{{p{\linewidth}}}]{
                \texttt{\tiny{[P1$\rangle$GM]}}
                \texttt{explanation: I am starting with a common five‑letter word that contains a mixture of common consonants (s, t, l) and a vowel (a, e) to get a broad sense of which letters are in the target word.  } \\
\texttt{guess: slate} \\
            }
        }
    }
    & & \\ \\

    \theutterance \stepcounter{utterance}  
    & & & \multicolumn{2}{p{0.3\linewidth}}{
        \cellcolor[rgb]{0.9,0.9,0.9}{
            \makecell[{{p{\linewidth}}}]{
                \texttt{\tiny{[GM$|$GM]}}
                \texttt{attempts = 1} \\
\texttt{target\_word = today} \\
\texttt{guess = slate} \\
\texttt{guess\_feedback = s<red> l<red> a<yellow> t<yellow> e<red>} \\
            }
        }
    }
    & & \\ \\

    \theutterance \stepcounter{utterance}  
    & & & \multicolumn{4}{p{0.6\linewidth}}{
        \cellcolor[rgb]{0.9,0.9,0.9}{
            \makecell[{{p{\linewidth}}}]{
                \texttt{\tiny{[P1$\langle$GM]}}
                \texttt{guess\_feedback: s<red> l<red> a<yellow> t<yellow> e<red>} \\
\\ 
\texttt{Provide your response only in this format.} \\
\texttt{explanation: details} \\
\texttt{guess: word} \\
            }
        }
    }
     \\ \\

    \theutterance \stepcounter{utterance}  
    & \multicolumn{4}{p{0.6\linewidth}}{
        \cellcolor[rgb]{0.561,0.737,0.561}{
            \makecell[{{p{\linewidth}}}]{
                \texttt{\tiny{[P1$\rangle$GM]}}
                \texttt{explanation: We know that s, l, and e are not in the word, while a and t are present but cannot be in the 3rd or 4th positions. By guessing "tango", we place a in the 2nd position and t in the 1st position—both allowed—and test three new letters (n, g, o) to learn more about the remaining spots.  } \\
\texttt{guess: tango} \\
            }
        }
    }
    & & \\ \\

    \theutterance \stepcounter{utterance}  
    & & & \multicolumn{2}{p{0.3\linewidth}}{
        \cellcolor[rgb]{0.9,0.9,0.9}{
            \makecell[{{p{\linewidth}}}]{
                \texttt{\tiny{[GM$|$GM]}}
                \texttt{attempts = 2} \\
\texttt{target\_word = today} \\
\texttt{guess = tango} \\
\texttt{guess\_feedback = t<green> a<yellow> n<red> g<red> o<yellow>} \\
            }
        }
    }
    & & \\ \\

    \theutterance \stepcounter{utterance}  
    & & & \multicolumn{4}{p{0.6\linewidth}}{
        \cellcolor[rgb]{0.9,0.9,0.9}{
            \makecell[{{p{\linewidth}}}]{
                \texttt{\tiny{[P1$\langle$GM]}}
                \texttt{guess\_feedback: t<green> a<yellow> n<red> g<red> o<yellow>} \\
\\ 
\texttt{Provide your response only in this format.} \\
\texttt{explanation: details} \\
\texttt{guess: word} \\
            }
        }
    }
     \\ \\

    \theutterance \stepcounter{utterance}  
    & \multicolumn{4}{p{0.6\linewidth}}{
        \cellcolor[rgb]{0.561,0.737,0.561}{
            \makecell[{{p{\linewidth}}}]{
                \texttt{\tiny{[P1$\rangle$GM]}}
                \texttt{explanation: The feedback tells us that the target word starts with t (green), contains a and o somewhere else (yellow), and does not contain n or g.  Since a cannot be in the second position and o cannot be in the fifth position, a viable 5‑letter word that satisfies all constraints is “toady.” This word places a in the third spot and o in the second spot, keeping both letters correct but misplaced as indicated.  } \\
\texttt{guess: toady} \\
            }
        }
    }
    & & \\ \\

    \theutterance \stepcounter{utterance}  
    & & & \multicolumn{2}{p{0.3\linewidth}}{
        \cellcolor[rgb]{0.9,0.9,0.9}{
            \makecell[{{p{\linewidth}}}]{
                \texttt{\tiny{[GM$|$GM]}}
                \texttt{attempts = 3} \\
\texttt{target\_word = today} \\
\texttt{guess = toady} \\
\texttt{guess\_feedback = t<green> o<green> a<yellow> d<yellow> y<green>} \\
            }
        }
    }
    & & \\ \\

    \theutterance \stepcounter{utterance}  
    & & & \multicolumn{4}{p{0.6\linewidth}}{
        \cellcolor[rgb]{0.9,0.9,0.9}{
            \makecell[{{p{\linewidth}}}]{
                \texttt{\tiny{[P1$\langle$GM]}}
                \texttt{guess\_feedback: t<green> o<green> a<yellow> d<yellow> y<green>} \\
\\ 
\texttt{Provide your response only in this format.} \\
\texttt{explanation: details} \\
\texttt{guess: word} \\
            }
        }
    }
     \\ \\

    \theutterance \stepcounter{utterance}  
    & \multicolumn{4}{p{0.6\linewidth}}{
        \cellcolor[rgb]{0.561,0.737,0.561}{
            \makecell[{{p{\linewidth}}}]{
                \texttt{\tiny{[P1$\rangle$GM]}}
                \texttt{explanation: The feedback shows that t, o, and y are in the correct positions (green).  The letters a and d are present (yellow) but are not in the current spots of the guess; thus, a must be at position 4 and d at position 3.  The only common five‑letter word fitting t‑o‑d‑a‑y is “today.”  } \\
\texttt{guess: today} \\
            }
        }
    }
    & & \\ \\

    \theutterance \stepcounter{utterance}  
    & & & \multicolumn{2}{p{0.3\linewidth}}{
        \cellcolor[rgb]{0.9,0.9,0.9}{
            \makecell[{{p{\linewidth}}}]{
                \texttt{\tiny{[GM$|$GM]}}
                \texttt{attempts = 4} \\
\texttt{target\_word = today} \\
\texttt{guess = today} \\
\texttt{guess\_feedback = t<green> o<green> d<green> a<green> y<green>} \\
            }
        }
    }
    & & \\ \\

    \theutterance \stepcounter{utterance}  
    & & & \multicolumn{2}{p{0.3\linewidth}}{
        \cellcolor[rgb]{0.9,0.9,0.9}{
            \makecell[{{p{\linewidth}}}]{
                \texttt{\tiny{[GM$|$GM]}}
                \texttt{game\_result = WIN} \\
            }
        }
    }
    & & \\ \\

\end{supertabular}
}

\end{document}
